\label{sub:IncreasingTheSizeOfExistingDPFs}
Another approach is hardcoding larger datatypes or scaling existing datatypes replacing the originals.
The user chooses her preferred data-size before compilation and uses this specified size during execution.
This will commit each duoram to a single size which, however, can be chosen freely before compilation.

The updated version comes in two flavors:
\begin{enumerate}
    \item \code{64ToUnlimited}\footnote{Public on https://github.com/FelixWeik/PrivateRandomAccessComputations/tree/64ToUnlimited} as found on the \hyperlink{https://github.com/FelixWeik/PrivateRandomAccessComputations/tree/64ToUnlimited}{branch of the same name in the author's repository}.
    It uses a boilerplate-like approach of hardcoding the same functionality for multiple hardcoded datasizes.
    The \hyperlink{https://github.com/FelixWeik/PrivateRandomAccessComputations/tree/64ToUnlimited}{updated implementation} supports 64, 128, and 256 bit values.
    \item \code{UnlimitedGMP}\footnote{Public on https://github.com/FelixWeik/PrivateRandomAccessComputations/tree/UnlimitedGMP} utilizes the \textit{GNU Multi Precision Arithmetic Library (GMP)}\footnote{More details can be found here: https://de.wikipedia.org/wiki/GNU\_Multiple\_Precision\_Arithmetic\_Library}.
    In general, it allows for implementing integral datatypes like integers with arbitrary precision while supporting general bit-operations.
    Thus, the updated code found on branch \hyperlink{https://github.com/FelixWeik/PrivateRandomAccessComputations/tree/UnlimitedGMP}{UnlimitedGMP} supports any size being a multiple of 128 bits (\ref{sec:GeneralLimitations}).
    Further details are found in \ref{subsub:gmp} and \ref{subsub:applyingGMP}.
\end{enumerate}

Both approaches build upon the existing implementation and follow the same basic approach.
Yet, they differ significantly taking a deeper look into their codebase.

\subsubsection*{Locating necessary Changes} \label{subsub:necessaryChanges}
Considering the source code's structure, the author chose a top-down approach for implementing this solution.
This includes locating the most general datatypes, distinguishing cases and assigning each case its special datatype.
The most optimal solution could trigger a trickle-down-effect, where general datatypes affect more specialized datatypes with the ultimate goal of changing a single datatype or even a single flag in order to adjust the whole process of storing and retreiving data of a given size.
Therefore, the required efforts of achieving the optimal solution for this approach extend the efforts of solely switching the datatype used by the RDPFs.

The topmost elements in  this datatype-top-to-bottom structure are macros (\ref{subsub:implementationTypeshpp}).
These are specified in \code{types.hpp}.
They specify the most general datatypes and thus their derivates, aliases.
\begin{itemize}
    \item Type-alias \code{value\_t} by itself is used 258 times throughout the code.
    \code{VALUE\_BITS} calls the according cases.
    Its most prominent usage is in specifying the size of the register types \code{RegAS} and \code{RegXS}, and the size of DPF-leafs.
    \item \code{ADDRESS\_MAX\_BITS} specifies the size of addresses.
    The derived type-alias \code{address\_t} is used 143 times throughout the base-implementation.
    It deceides the type of indexing into the database.
    As the database is of size $2^\text{INPUT\_BITS}$ at least \code{INPUT\_BITS} many bits are needed for formatting adresses correctly.
\end{itemize}

\subsubsection*{Bitmasks and Bitutils}\label{subsub:bitmasksAndBitutils}
\code{bitutils.hpp} provides bitmasks and important functions for bit operations.
These operations include \code{xor\_if()} and \code{get\_lsb()} which are used multiple times throughout the code.
\code{xor\_if()} will execute a bitwise xor operation if a flag-bit is set, \code{get\_lsb()} retrieves the least significant bit of a given \code{\_\_m128i} value.

Bitmasks like \code{bool128\_mask} and \code{if128\_mask} are utilised when generating and traversing DPF-trees, during bit-operations, or while blinding shares.
In short, they increase bit-operations' performance by providing predefined values that will be used frequently during execution.

Since these bit-operations are executed primarily on \code{\_\_m128i} values, all function signatures and bitmasks are predefined for this specific datatype.
If one wants to scale their size arbitrarily, two solutions exist:
\begin{enumerate}
    \item Cut larger values into 128-bit blocks:
    This requires the input-type to be a multiple of 128-bit, which is not always the case.
    Also, it requires each function call to cut the input into $N$ blocks, call the given function $n$ times and finally concatenate each result to receive the final result.
    This too scales badly.
    \item Using GMP (\ref{sub:approachGMP}):
    By using GMP, one can create datatypes of arbitrary precision.
    Also, bit-operations are predefined for any \code{mpz} value.
    Thus, the specialized bit-operations defined in the code can be implemented to utilize the more general GMP bit-operations.
    Bitmasks can be left untouched, since the least significant bit or bit-inverses are not size-dependent.
    By using this library for bit-operations while leaving the bitmasks as they are, one achieves an arbitrarily scaled set of bit-utils.
    Therefore, the author further pursues this approach.
\end{enumerate}

The approach on GMP will be further discussed in the following section \ref{sub:approachGMP}.

\subsubsection*{Adjusting Compilation}\label{subsub:adjustingCompilation}
Both approaches require a flag bit to decide on which datatype or size to use during execution.
A new compilation-flag \code{-DVALUE\_BITS} is set in the \code{MAKEFILE} and represents the length of data one wants to store and retrieve.
It directly affects \code{VALUE\_BITS} and places the given value into it.
\code{VALUE\_BITS} then triggers the desired trickle-down affect which further changes macros and datatypes.
If nothing else is specified, \code{-DVALUE\_BITS} is set to $64$ by default which is similar to the base-implementation of Vadapalli et al.

\subsubsection*{Introducing new Tests}
%//TODO hier gehts weiter lieber Felix

\subsubsection*{General Discussion}